\section{Mechanization and Trusted Base}
\label{sec:artifact}

The artifact contains 27,230 physical lines of Rocq source.
Since our Rocq development is heavily AI-assisted,
this line count does not accurately reflect our engineering effort.
Appendix~\ref{app:ai-assistance} describes the roles of the authors and AI tools in
the development and manuscript preparation.
%the line count measures the scope of the checked development rather than human effort.
The overview in \Cref{fig:theory-map} separates the application-specific
security argument from the reusable probability and program-logic results.

The mechanization contains a few layers.  The scheme layer
defines encryption games and the noise-flooded transform.  The
probability layer develops facts about countable discrete distributions,
without reference to programs.  The logic layer lifts those facts to SSProve
code.  The compiler layer implements the trace transformation and derives the
generic adaptive replacement rule.  Finally, the security
layer instantiates the reusable logic and compiler with the decryption
invariant, constructs the IND-CPA reduction, and composes the exact and
quantitative game transitions.
The top-level result is a functor over precisely the scheme, chart,
probability-one correctness, IND-CPA security, and flooding-parameter
interfaces described in \Cref{sec:construction}; probability lemmas and
program rules are proved dependencies.

\subsection{Formalization of judgments and rules}

We follow SSProve's convention and engineering framework: pure computation is
embedded in Rocq, while probabilistic and stateful effects are represented by
SSProve's free-monadic program syntax, packages, heaps, and linking operations
\cite{ssprove}.  The new logic does not extend Rocq with a primitive judgment
or a collection of trusted proof rules.
Judgments are formalized as ordinary Rocq propositions over the denotation \texttt{Pr\_code},
and inference rules are lemmas which prove such judgments.

For example, the exact definition of the additive-error judgment is shown
below.  Its types range over SSProve programs and heaps; \texttt{complete}
adds an explicit outcome for missing subdistribution mass.

\begingroup
\setmonofont{FreeMono}[Scale=MatchLowercase]
\inputminted[
  fontsize=\scriptsize,
  breaklines,
  linenos,
  firstnumber=58,
  firstline=58,
  lastline=70
]{coq}{rocq-excerpts/theories/ProgramLogics/Ae.v}
\endgroup

Thus, for every related pair of inputs and initial heaps, the proposition
requires a coupling of the completed output-and-heap distributions in which
the postcondition holds with probability at least \(1-\varepsilon\).
The Pythagorean judgment is encoded in the same style.  For each related input
pair it existentially quantifies two full transcript distributions, requires
their final marginals to equal the two completed program denotations, and
requires their supports to satisfy the postcondition.  Its predicate
\texttt{pythDist} additionally imposes full mass and the coordinatewise
conditional-KL bounds from \Cref{def:pyth-judgment}.
To place every result type in one transcript space, the Rocq definition
injectively encodes typed output/heap pairs in \texttt{nat * heap} and
represents \(\bot\) by \texttt{None}.

The rules in \Cref{fig:logic-rules} are proved implications between these
definitions.  For instance, the source-level statement of the final
KL-to-additive-error bridge is:

\begingroup
\setmonofont{FreeMono}[Scale=MatchLowercase]
\inputminted[
  fontsize=\scriptsize,
  breaklines,
  linenos,
  firstnumber=530,
  firstline=530,
  lastline=543
]{coq}{rocq-excerpts/theories/ProgramLogics/Pyth.v}
\endgroup

Here \texttt{pythagorean\_tv\_bound s} is
\(\sqrt{\sum_i s_i/2}\).  Its proof applies the verified probability theorem
to the transcript witnesses, projects to their final marginals, and invokes
the completed-distribution coupling theorem.  The sampling, sequencing,
mixed Hoare/Pythagorean, and compiler rules have the same status: each is a
Rocq lemma proved from the program semantics.  Consequently, a Rocq kernel check of
the security theorem also checks the soundness chain from the distribution
lemmas through the program judgments.

\subsection{Compiler implementation and correctness}

\Cref{fig:compiler-code} shows the traversal over SSProve's \texttt{raw\_code}
syntax.  Its result pairs the trace with a status: \texttt{inl} carries the
encoded input of a selected call, whereas \texttt{inr} carries the final output
of a program that terminated first.  A nonselected call is executed, its
result is appended to the trace, and traversal continues; heap operations and
samples are handled analogously.

\begin{figure}[t]
\begingroup
\setmonofont{FreeMono}[Scale=MatchLowercase]
\inputminted[
  fontsize=\scriptsize,
  breaklines,
  linenos,
  firstnumber=167,
  firstline=167,
  lastline=178
]{coq}{rocq-excerpts/theories/NextMessage/Trace.v}
\endgroup
\caption{Scanning for the next call to \texttt{fn}: \texttt{inl} reports its
encoded input, while \texttt{inr} reports program termination.}
\label{fig:compiler-code}
\end{figure}

The outer driver answers a discovered call, records the answer, and repeats.
After the requested calls, trace replay reconstructs the remaining code, so
the complete transformation has the original result type.  The same-package
correctness result is \Cref{thm:compile-same}, whose Rocq theorem is
\texttt{compile\_calls\_correct}.  It proves equality of program semantics by
induction over traversal and trace replay.

The separate quantitative result is \Cref{thm:compile-replace}.  Its call
induction lifts the one-call Pythagorean judgment through the trace
transformation and accumulates \(q\) copies of its KL-cost tuple;
\texttt{compileRule} then performs the single final Micciancio--Walter
conversion.  The \textsc{Compile} rule in
\Cref{fig:security-rule-inventory} records its complete security-facing
premises.  Thus both semantic compiler correctness and adaptive quantitative
replacement are checked consequences of SSProve syntax and semantics.

\subsection{Probability formalization}

The probability development works over MathComp's countable discrete
distributions, rather than restricting the Gaussian argument to finite
support.  It defines KL finiteness as absolute continuity together with
summability of the KL integrand, proves Pinsker's inequality, and proves the
conditional-coordinate chain bound used in
\Cref{thm:pyth-probability}.  It also constructs completion and maximal
couplings, so that a distribution-level total-variation bound yields the
program-level proposition above.

The distributions used here, notably the discrete Gaussians, have countably
infinite support.  Consequently, the KL chain rule involves infinite series.
A paper proof may leave their convergence implicit; the formalization must
establish it before decomposing transcript KL into conditional KL costs.  It
also treats histories of probability zero explicitly.  These results yield the
chain bound used in the square-root composition argument.

For the one-call bound, the development also proves that the discrete Gaussian
is normalized and has full support, that its normalizer is invariant under
integer shifts, and that its centered first moment is zero.  These facts yield
the exact equal-variance KL identity in \Cref{eq:dg-kl} and discharge the
Gaussian sampling rule.  Further theorem statements and proof structure
appear in \Cref{sec:verified-kl-analysis}.

\subsection{The checked theorem}

At the application boundary, the functor defines the concrete reduction
\(\cB_{\cA,q}\) and makes the assumed IND-CPA bound and the statistical loss
explicit:

\begingroup
\setmonofont{FreeMono}[Scale=MatchLowercase]
\inputminted[
  fontsize=\scriptsize,
  breaklines,
  linenos,
  firstnumber=531,
  firstline=531,
  lastline=534
]{coq}{rocq-excerpts/theories/Security/NoiseFloodingSecurity/Final.v}
\endgroup

After proving that the constructed reduction has the IND-CPA adversary
interface and deriving the completed-output inequality, the development
exports the following theorem.

\begingroup
\setmonofont{FreeMono}[Scale=MatchLowercase]
\inputminted[
  fontsize=\scriptsize,
  breaklines,
  linenos,
  firstnumber=1954,
  firstline=1954,
  lastline=1957
]{coq}{rocq-excerpts/theories/Security/NoiseFloodingSecurity/Final.v}
\endgroup

The premise says that \(A\) is a well-typed package with the required IND-CPAD
adversary interface.  Unfolding \texttt{security\_bound} and
\texttt{security\_loss} gives exactly \Cref{thm:main}; in particular, the
right-hand term is the assumed IND-CPA winning-probability bound for the
constructed reduction plus \(\sqrt{qn}/(2\gamma)\).  This source statement also
makes clear where the theorem boundary lies: scheme structure, charts,
probability-one correctness, IND-CPA security, and positivity of \(\gamma\)
enter through the enclosing functor, whereas reduction validity, the
discrete-Gaussian calculation, compiler correctness, and logic soundness have
already been proved.

\subsection{Reproducibility}

The complete artifact is available at
\url{https://github.com/ethanlee515/mending}.
The complete artifact checks with Rocq 9.0.1, MathComp 2.6.0, MathComp Analysis and
experimental reals 1.18.0, and \texttt{rocq-ssprove} from its \texttt{main}
branch. The SSProve revision used for this upgrade is
\texttt{9a37fe0ab79d}; the artifact README records its full hash and commands
for selecting that snapshot. MathComp's built-in ring tactic replaces the
separate algebra-tactics package. The branch pin can advance when refreshed,
so subsequent dependency updates require a complete rebuild and assumption
audit. The manuscript's
Rocq listings are refreshed from the checked sources during its build, so the
displayed propositions and theorem statements do not silently diverge from
the artifact.

\subsection{Trusted computing base}
\label{sec:tcb}

The trusted base includes the Rocq kernel and standard library; MathComp,
MathComp Analysis, and their real-number/classical infrastructure; and
SSProve's program, heap, package, linking, and subdistribution semantics.
The scheme, chart, correctness, parameter, and IND-CPA interfaces from
\Cref{sec:construction} are theorem parameters and must be discharged by a
concrete instantiation.

The new KL, Gaussian, relational, compiler, and game-hop results contain no
local \texttt{Admitted}. Early versions of the artifact did inherit an
unproved summation assumption through SSProve's distribution semantics.
MathComp Analysis already contained statements of several summation lemmas,
but their proofs had been left admitted. One of these lemmas, which justifies
exchanging the order of countable summation, appeared in the assumption audit
of our top-level security theorem.

We proved these lemmas and contributed the proofs upstream, where they were
accepted and merged into MathComp Analysis.
SSProve now supports the updated library, so the artifact uses the upstream
proofs directly. After rebuilding the complete artifact with these
dependencies, we reran the top-level assumption audit and confirmed that it
no longer reports an admitted summation lemma.

The remaining foundational assumptions reported by \texttt{Print Assumptions}
are MathComp's propositional extensionality, dependent functional
extensionality, and constructive indefinite description. Within the functor,
the audit also lists SSProve's abstract real type \(R\), the scheme, chart,
correctness, and IND-CPA inputs, and positivity of the discrete-Gaussian width
multiplier \(\gamma\). A concrete scheme must discharge the application
interfaces. The audit does not report any program-logic rule, compiler
correctness theorem, or discrete-Gaussian KL identity: these are verified
inside the development.

Selected definitions of the scheme interface, games, transform,
and reduction remain in \Cref{app:formal-interface-excerpts}.  Together with
the source-level theorem above, they expose the small specification surface on
which a cryptographic audit should concentrate.
\Cref{tab:formal-excerpt-size} quantifies the selected appendix excerpts;
the Rocq kernel checks that the
remaining proof terms connect that surface by the argument in
\Cref{sec:verified-reduction}.  As in SSProve and EasyCrypt, running time is not
part of the program semantics and must be inspected separately
\cite{easycrypt,ssprove}.  Here the
reduction runs the adversary once, forwards its oracle calls, and adds table
operations and discrete-Gaussian sampling.
